\documentclass[../main.tex]{subfiles}

\begin{document}

\section{Minimos Quadrados}

\subsection{Metodo implementado e objetivo numerico}

O metodo dos minimos quadrados foi usado para ajustar funcoes simples a dados
experimentais. A ideia central e escolher os coeficientes de um modelo de forma
que a soma dos quadrados dos residuos seja minima. Para pontos
\((x_i,y_i)\), o residuo e
\[
    r_i = y_i - p(x_i),
\]
e o criterio minimizado e
\[
    E = \sum_{i=1}^{n} r_i^2.
\]

No projeto foram implementados tres casos principais. O primeiro e a regressao
linear, em que \(p(x)=a_0+a_1x\). O segundo e o ajuste polinomial discreto, no
qual o grau pode ser escolhido livremente. O terceiro e o modelo de potencia,
\(y=ax^b\), resolvido por linearizacao logaritmica.

\subsection{Implementacao em Python}

Os codigos principais usados nesta parte sao:

\begin{itemize}
    \item \texttt{metodos/regressao\_linear\_mmq.py}, para a regressao linear,
    coeficiente de correlacao, \(R^2\) e residuos;
    \item \texttt{metodos/mmq\_discreto\_polinomial.py}, para ajuste de
    polinomios de qualquer grau por equacoes normais;
    \item \texttt{metodos/sistemas\_lineares.py}, com eliminacao de Gauss com
    pivoteamento parcial;
    \item \texttt{metodos/polinomios.py}, com avaliacao e formatacao de
    polinomios;
    \item \texttt{testes\_livro\_neide.py}, que organiza os exercicios do livro,
    gera CSVs detalhados e produz os graficos.
\end{itemize}

A leitura dos dados e feita a partir de arquivos TXT em UTF-8. Cada linha de
entrada contem os valores numericos do problema, e linhas iniciadas por
\texttt{\#} sao ignoradas. Os resultados sao gravados em CSV tambem em UTF-8,
usando a funcao auxiliar \texttt{salvar\_csv}.

Para o ajuste polinomial de grau \(m\), o sistema normal implementado e:
\[
    \sum_{j=0}^{m} c_j \sum_{i=1}^{n} x_i^{j+k}
    =
    \sum_{i=1}^{n} y_i x_i^k,
    \quad k=0,1,\ldots,m.
\]
Esse sistema e montado manualmente com somatorios e resolvido pela eliminacao
de Gauss do proprio projeto.

\subsection{Exercicio 8.1: acidentes em veiculos motorizados}

\paragraph{Enunciado e objetivo numerico}
O exercicio 8.1 apresenta dados sobre acidentes em veiculos motorizados no
Brasil entre 1980 e 1998. O item (a) pede uma regressao linear do numero de
acidentes no tempo, e o item (b) pede uma regressao quadratica para a taxa de
acidentes por 10000 veiculos. Em ambos os casos, o objetivo e prever o valor
para o ano 2000.

\paragraph{Metodo empregado}
No item (a), foi usada regressao linear por minimos quadrados. No item (b), foi
usado minimos quadrados polinomial de grau 2. Seguindo a orientacao do
enunciado, a variavel independente foi deslocada:
\[
    x = ano - 1980.
\]
Isso evita operar diretamente com numeros grandes como 1980, 1985 e 1998.

\paragraph{Desenvolvimento matematico}
Para os acidentes em milhares, o modelo usado foi
\[
    y = a_0 + a_1x.
\]
Para a taxa de acidentes por 10000 veiculos, o modelo usado foi
\[
    y = a_0 + a_1x + a_2x^2.
\]
Em ambos os casos, os coeficientes foram obtidos pelas equacoes normais do
metodo dos minimos quadrados.

\paragraph{Implementacao}
Os dados foram lidos do TXT do exercicio. O script principal de testes separa
as colunas da tabela em dois conjuntos de pontos: acidentes em milhares e
acidentes por 10000 veiculos. Depois chama a rotina de ajuste polinomial. Os
CSVs detalhados foram salvos separadamente, um para o modelo linear dos
acidentes e outro para o modelo quadratico da taxa.

\paragraph{Resultados}
Para o numero de acidentes, o ajuste obtido foi:
\[
    y = 8.02159 + 0.362485x.
\]
Para \(x=20\), correspondente ao ano 2000:
\[
    y(20)=15.2713.
\]
Logo, a previsao foi de aproximadamente \(15.2713\) milhares de acidentes.

Para a taxa de acidentes por 10000 veiculos, o polinomio ajustado foi:
\[
    y = 1.6985 - 0.0368563x + 0.00115007x^2.
\]
No ano 2000:
\[
    y(20)=1.42140.
\]

\[
\begin{array}{c|c|c}
\text{modelo} & \text{soma dos residuos ao quadrado} & R^2 \\
\hline
\text{acidentes, grau 1} & 2.123973 & 0.938071 \\
\text{taxa, grau 2} & 0.004692 & 0.939921
\end{array}
\]

\begin{figure}[H]
    \centering
    \includegraphics[width=0.82\textwidth]{ex8_1_acidentes_linear.png}
    \caption{Regressao linear do numero de acidentes no exercicio 8.1.}
\end{figure}

\begin{figure}[H]
    \centering
    \includegraphics[width=0.82\textwidth]{ex8_1_taxa_quadratica.png}
    \caption{Regressao quadratica da taxa de acidentes no exercicio 8.1.}
\end{figure}

\paragraph{Interpretacao}
O grafico da regressao linear mostra uma tendencia geral de crescimento do
numero de acidentes, embora alguns pontos intermediarios se afastem da reta. O
valor de \(R^2\) proximo de \(0.94\) indica que a reta explica boa parte da
variacao observada. Para a taxa por 10000 veiculos, a parabola captura uma
queda inicial e uma leve retomada ao final, o que justifica o uso do termo
quadratico.

\subsection{Exercicio 8.5: coeficiente de placa de orificio}

\paragraph{Enunciado e objetivo numerico}
O exercicio 8.5 fornece valores experimentais do coeficiente \(C\) em funcao da
razao \(x=A/A_1\). O objetivo e ajustar os pontos por
\[
    C(x)=a_0+a_1x+a_2x^2.
\]

\paragraph{Metodo empregado}
Foi usado minimos quadrados polinomial discreto de grau 2, pois ha mais pontos
experimentais do que coeficientes. Assim, em vez de interpolar todos os dados, o
programa encontra a parabola que minimiza a soma dos quadrados dos residuos.

\paragraph{Desenvolvimento matematico}
O sistema normal montado pelo codigo foi:
\[
\begin{bmatrix}
10 & 5.5 & 3.85\\
5.5 & 3.85 & 3.025\\
3.85 & 3.025 & 2.5333
\end{bmatrix}
\begin{bmatrix}
a_0\\a_1\\a_2
\end{bmatrix}
=
\begin{bmatrix}
7.4\\4.391\\3.2319
\end{bmatrix}.
\]
Resolvendo esse sistema por eliminacao de Gauss, obtem-se os coeficientes do
polinomio ajustado.

\paragraph{Implementacao}
Os dados foram lidos do TXT do exercicio. A funcao de ajuste polinomial monta
as equacoes normais e chama o resolvedor de Gauss. O sistema normal foi salvo
em um CSV proprio, e a tabela completa de comparacao foi salva no CSV principal
do exercicio.

\paragraph{Resultados}
O polinomio obtido foi:
\[
    C(x)=0.650167 - 0.231742x + 0.564394x^2.
\]

\[
\begin{array}{c|c|c|c}
x & C\text{ observado} & C\text{ ajustado} & \text{residuo} \\
\hline
0.10 & 0.62 & 0.632636 & -0.012636 \\
0.50 & 0.68 & 0.675394 & 0.004606 \\
0.80 & 0.81 & 0.825985 & -0.015985 \\
1.00 & 1.00 & 0.982818 & 0.017182
\end{array}
\]

\begin{figure}[H]
    \centering
    \includegraphics[width=0.82\textwidth]{ex8_5_orificio.png}
    \caption{Ajuste quadratico para o coeficiente \(C\) no exercicio 8.5.}
\end{figure}

\paragraph{Interpretacao}
O ajuste acompanha a tendencia crescente dos dados. Os maiores residuos em
modulo aparecem proximos das extremidades, principalmente em \(x=1\), onde o
valor tabelado e \(1.00\) e o valor ajustado e \(0.982818\). Ainda assim, os
residuos sao pequenos para a escala do problema, e o grafico mostra uma curva
coerente com os pontos experimentais.

\subsection{Exercicio 8.11: modelo de potencia para o PNB}

\paragraph{Enunciado e objetivo numerico}
O exercicio 8.11 fornece o Produto Nacional Bruto em dolares correntes e em
dolares constantes. O enunciado pede que ambas as series sejam aproximadas por
uma funcao do tipo:
\[
    y = ax^b.
\]
Depois, o modelo deve ser usado para prever os valores no ano 2000.

\paragraph{Metodo empregado}
O modelo de potencia nao e linear nos parametros originais. Por isso, foi usada
a transformacao logaritmica:
\[
    \ln y = \ln a + b\ln x.
\]
Assim, o problema vira uma regressao linear entre \(X=\ln x\) e \(Y=\ln y\).
Como \(x\) precisa ser positivo, usei \(x=ano-1979\).

\paragraph{Implementacao}
A funcao de regressao por potencia transforma os pontos por logaritmo, chama o
ajuste linear e depois recupera \(a=e^{c_0}\). Os dados de entrada estao no TXT
do exercicio. O CSV resumido guarda os parametros dos dois modelos, e a tabela
detalhada registra valores observados, valores estimados e residuos.

\paragraph{Resultados}
Para o PNB em dolares correntes:
\[
    y = 225.107890x^{0.425689},
    \quad y(21)=822.7053.
\]
Para o PNB em dolares constantes:
\[
    y = 329.451061x^{0.236483},
    \quad y(21)=676.8197.
\]
A razao entre as previsoes foi:
\[
    \frac{822.7053}{676.8197}=1.21555.
\]

\begin{figure}[H]
    \centering
    \includegraphics[width=0.82\textwidth]{ex8_11_pnb_corrente_constante.png}
    \caption{Ajuste por modelo de potencia para os PNBs do exercicio 8.11.}
\end{figure}

\paragraph{Interpretacao}
O modelo indica que o PNB em dolares correntes cresce mais rapidamente do que o
PNB em dolares constantes. A razao \(1.21555\) sugere que, para 2000, o valor em
dolares correntes seria cerca de \(21.55\%\) maior que o valor em dolares
constantes. Numericamente, essa diferenca representa o efeito acumulado da
inflacao no modelo ajustado.

\subsection{Validacao dos resultados de minimos quadrados}

Os graficos confirmam que os ajustes seguem a tendencia principal dos dados. Em
8.1, os dois modelos apresentaram \(R^2\) proximo de \(0.94\), indicando boa
capacidade explicativa para uma tabela pequena. Em 8.5, os residuos sao
pequenos e alternam sinais, comportamento esperado em um ajuste por minimos
quadrados. Em 8.11, a transformacao logaritmica permitiu tratar o modelo de
potencia com as mesmas ferramentas da regressao linear, sem usar rotinas prontas
de ajuste.

\end{document}
